\chapter{Conclusiones}
\label{sec.ej2:conclusiones}

El análisis y la simulación del módem de FM mediante el método de demodulación por derivación permitieron comprender
de manera práctica y matemática el funcionamiento de los sistemas de modulación angular. Se validaron con éxito los
conceptos de desviación de frecuencia, índice de modulación y ancho de banda de Carson.

A partir del desarrollo y los resultados obtenidos, se pueden extraer las siguientes conclusiones clave:
\begin{itemize}
    \item El método de demodulación por derivación constituye una alternativa simple y de bajo costo para recuperar
          señales de FM. Su principio se basa en convertir las variaciones de frecuencia en variaciones de amplitud
          mediante un bloque derivador, cuya ganancia es directamente proporcional a la frecuencia instantánea.
    \item Se demostró analítica y experimentalmente la existencia de un límite estricto para el índice de modulación,
          dado por la relación $m_f \le f_c / f_m$. Para las frecuencias de simulación ensayadas, el índice de
          modulación máximo sin distorsión es $m_{f,\text{máx}} = 10$. Superar dicho límite provoca sobremodulación
          en la componente AM resultante de la derivación, lo que se traduce en inversiones de fase de la portadora
          y una severa distorsión armónica en la señal de salida demodulada.
    \item La distribución de bandas laterales obtenida en el analizador de espectro concuerda con las funciones de
          Bessel para un índice $m_f = 5$. Asimismo, el ancho de banda efectivo cumple estrictamente con el valor
          de $12\kHz$ estimado mediante la regla de Carson.
    \item A pesar de su sencillez, el discriminador por derivación presenta dos limitaciones críticas que restringen
          su uso en aplicaciones reales de telecomunicaciones si no se incorporan etapas adicionales:
          \begin{enumerate}
              \item \emph{Sensibilidad al ruido de amplitud}: Debido a que la etapa de detección final responde a la
                    envolvente de amplitud de la señal, cualquier ruido o interferencia que altere la amplitud de la
                    portadora se trasladará directamente a la señal demodulada. En la práctica, esto requiere colocar
                    una etapa limitadora de amplitud antes del derivador.
              \item \emph{Amplificación del ruido de alta frecuencia}: La respuesta en frecuencia de un bloque
                    derivador es linealmente creciente con la frecuencia ($H(\omega) \propto j\omega$). Esto significa
                    que el derivador actúa como un filtro pasa altos, amplificando en mayor medida los ruidos de alta
                    frecuencia que se encuentran en el canal, empeorando la relación señal-ruido del receptor.
          \end{enumerate}
\end{itemize}

En síntesis, la simulación realizada permitió contrastar y validar la teoría matemática de FM, confirmando la
factibilidad del demodulador por derivación bajo condiciones de diseño controladas y delimitando con precisión sus
umbrales operativos de no distorsión.
